\documentclass[a4paper,14pt]{extarticle}
\usepackage[utf8]{inputenc}
\usepackage[T2A]{fontenc}
\usepackage[english,russian]{babel}
\usepackage{geometry}
\geometry{left=25mm, right=15mm, top=20mm, bottom=20mm}
\usepackage{indentfirst}
\usepackage{graphicx}
\usepackage{float}
\usepackage{hyperref}

\begin{document}

\begin{titlepage}
\begin{center}
\vspace*{2cm}
{\Large \textbf{ТЕХНИЧЕСКИЙ ПРОЕКТ}} \\
\vspace{0.5cm}
{\large на программу} \\
\vspace{0.5cm}
{\huge \textbf{«Motion Visualizer»}} \\
\vspace{0.5cm}
\vfill
г. Ростов-на-Дону, 2026
\end{center}
\end{titlepage}

\tableofcontents
\newpage

\section{Архитектура программы}

\subsection{Общая структура}
Программа состоит из следующих модулей:
\begin{itemize}
    \item \textbf{main.py} — точка входа, главное окно (PySide6)
    \item \textbf{glwidget.py} — 3D-визуализация (OpenGL)
    \item \textbf{motion\_core.so} — C++ модуль расчёта движения
    \item \textbf{db\_client.py} — работа с SQLite
\end{itemize}

\subsection{Схема взаимодействия модулей}
\begin{figure}[H]
\centering
\framebox(400,100){GUI → C++ → OpenGL → SQLite}
\caption{Схема потоков данных}
\end{figure}

\section{Алгоритмы}

\subsection{Интегратор движения}
Используется метод Эйлера:
\[
x_{n+1} = x_n + v_x \cdot \Delta t
\]
\[
y_{n+1} = y_n + v_y \cdot \Delta t
\]

\subsection{Обработка событий мыши}
При перетаскивании объекта вычисляется смещение курсора и обновляется позиция:
\[
\Delta x = x_{current} - x_{last}
\]
\[
\Delta y = y_{current} - y_{last}
\]

\section{Стек технологий}

\begin{table}[H]
\centering
\begin{tabular}{|p{5cm}|p{8cm}|}
\hline
\textbf{Компонент} & \textbf{Технология} \\
\hline
GUI & PySide6 (Qt6) \\
\hline
3D-графика & PyOpenGL + QtOpenGLWidgets \\
\hline
Вычислительное ядро & C++20 + pybind11 \\
\hline
База данных & SQLite + SQLAlchemy \\
\hline
Сборка Python & uv \\
\hline
Сборка C++ & CMake \\
\hline
\end{tabular}
\caption{Стек технологий}
\end{table}

\end{document}